\documentclass[11pt]{article}
\usepackage[a4paper,margin=2.5cm]{geometry}
\usepackage{amsmath,amssymb,amsthm}
\usepackage[british]{babel}
\usepackage{microtype}
\setlength{\emergencystretch}{3em}
\newtheorem{proposition}{Proposition}
\newtheorem{lemma}[proposition]{Lemma}
\theoremstyle{remark}
\newtheorem{remark}[proposition]{Remark}
\newcommand{\E}{\mathbb{E}}
\newcommand{\Q}{\mathbb{Q}}
\renewcommand{\P}{\mathbb{P}}

\title{Crash insurance and the G10 carry premium: theory notes}
\author{Amanjeet Singh}
\pdfinfo{
  /Author (Amanjeet Singh)
  /Title (Crash insurance and the G10 carry premium: theory notes)
  /Subject (Proofs of theory results R1 to R10 except R9, which is cited with its conditions)
  /Keywords (carry trade, crash risk, FX options, option-implied moments, partial identification)
  /Creator (LaTeX)
}
\date{Version of 3 October 2026}

\begin{document}
\maketitle

\noindent These notes prove every result listed in \texttt{theory/README.md} except R9 (inference), which is cited there with its conditions and is not proved here. Results that are standard are marked as such and attributed, and where a published proof follows the same route as the one written here, it is cited. The parts I claim as my own are stated at the start of each section. Throughout, $\Phi$ and $\varphi$ are the standard normal distribution and density functions, $s=\sigma\sqrt{\tau}$, and for forward $F$ and strike $K$
\[
d_\pm(K)=\frac{\ln(F/K)\pm s^2/2}{s}.
\]
The identity $K\varphi(d_-)=F\varphi(d_+)$ is used repeatedly.

\section*{R1. Domestic and foreign risk-neutral measures}

This is the change of numeraire of Geman, El Karoui and Rochet (1995); the statement and proof are included for completeness. Girsanov's theorem and It\^o's formula are standard results, stated for example in Shreve (2004).

Let $(\Omega,\mathcal F,(\mathcal F_t)_{t\in[0,T]},\Q^d)$ satisfy the usual conditions, let $B^d,B^f$ be strictly positive adapted money-market accounts with $B^d_0=B^f_0=1$, and let $S>0$ be adapted and c\`adl\`ag. Assume $Z_t=S_tB^f_t/(S_0B^d_t)$ is a true $\Q^d$-martingale and define $d\Q^f/d\Q^d=Z_T$ on $\mathcal F_T$.

\begin{proposition}
(a) $\Q^f$ is a probability measure equivalent to $\Q^d$. (b) For $\mathcal F_T$-measurable $X\ge0$ or $X\in L^1(\Q^f)$, $\E^{\Q^f}[X\mid\mathcal F_t]=\E^{\Q^d}[Z_TX\mid\mathcal F_t]/Z_t$. (c) An adapted process $Y$ with $Y_t\in L^1(\Q^f)$ is a $\Q^f$-martingale if and only if $ZY$ is a $\Q^d$-martingale; in particular a foreign price $V^f$ satisfies: $V^f/B^f$ is a $\Q^f$-martingale if and only if $SV^f/B^d=S_0Z\,(V^f/B^f)$ is a $\Q^d$-martingale. (d) If $dS_t/S_t=(r_d-r_f)\,dt+\sigma\,dW_t$ under $\Q^d$ with constant coefficients, then $W^f_t=W_t-\sigma t$ is a $\Q^f$-Brownian motion, $d(1/S_t)=(1/S_t)[(r_f-r_d)\,dt-\sigma\,dW^f_t]$, and $\E^{\Q^d}S_T=F$, $\E^{\Q^f}(1/S_T)=1/F$ with $F=S_0e^{(r_d-r_f)T}$.
\end{proposition}

\begin{proof}
(a) $Z_T>0$ and $\E^{\Q^d}Z_T=Z_0=1$, so $\Q^f$ is a probability measure; positivity of $Z_T$ gives equivalence. (b) For $\mathcal F_t$-measurable $Y\ge0$, the tower property and the martingale property give $\E^{\Q^f}Y=\E^{\Q^d}[Z_TY]=\E^{\Q^d}[Z_tY]$. For $A\in\mathcal F_t$,
\[
\E^{\Q^f}[X\mathbf 1_A]=\E^{\Q^d}\bigl[\E^{\Q^d}[Z_TX\mid\mathcal F_t]\mathbf 1_A\bigr]=\E^{\Q^f}\Bigl[\frac{\E^{\Q^d}[Z_TX\mid\mathcal F_t]}{Z_t}\mathbf 1_A\Bigr],
\]
which characterises the conditional expectation. (c) Apply (b) on $[s,t]$: $\E^{\Q^f}[Y_t\mid\mathcal F_s]=\E^{\Q^d}[Z_tY_t\mid\mathcal F_s]/Z_s$, so $Y$ is a $\Q^f$-martingale exactly when $\E^{\Q^d}[Z_tY_t\mid\mathcal F_s]=Z_sY_s$. (d) Here $Z_t=\exp(\sigma W_t-\sigma^2t/2)$, so Girsanov's theorem gives the first claim, and It\^o's formula for $1/S$ gives the dynamics. The expectations follow from the lognormal law under each measure.
\end{proof}

\begin{remark}
Jensen's inequality gives $\E^{\Q^d}(1/S_T)>1/\E^{\Q^d}S_T=1/F$. There is no contradiction with (d): the foreign investor's forward price is an expectation under $\Q^f$, not $\Q^d$ (Siegel's paradox). If $Z$ is only a strict local martingale, then $\E Z_T<1$ and $Z_T$ does not define an equivalent probability measure. Carr, Fisher and Ruf (2014) construct the foreign measure in that case as a F\"ollmer measure that is not equivalent to $\Q^d$, and restore the domestic--foreign symmetry with a modified pricing operator.
\end{remark}

\section*{R2. Arbitrage-free call prices}

The link between the second strike derivative of call prices and state prices is due to Breeden and Litzenberger (1978). The characterisation below is a known result: F\"ollmer and Schied (2004, Lemma 7.23) show that call prices consistent with a linear, positive pricing rule are the call prices of a unique probability measure with the forward as its mean, by the same argument through the right derivative of the call-price function and Fubini's theorem. I state it in the form used by the arbitrage checks and write out the proof for completeness.

\begin{proposition}
Let $c:[0,\infty)\to\mathbb{R}$ and $F>0$. There is a random variable $X\ge 0$ with $\E X=F$ and $c(K)=\E(X-K)^+$ for all $K\ge0$ if and only if (i) $c$ is convex, (ii) $c(0)=F$, (iii) $c'_+(0)\ge -1$ and (iv) $c(K)\to0$ as $K\to\infty$. The law of $X$ is then unique, with distribution function $G=1+c'_+$ on $[0,\infty)$ and $\P(X=0)=1+c'_+(0)$.
\end{proposition}

\begin{proof}
\emph{Necessity.} Convexity of $c$ follows from convexity of $K\mapsto(x-K)^+$ and linearity of expectation, and $c(0)=\E X=F$. Since $0\le(X-K)^+\le X\in L^1$ and $(X-K)^+\to0$ pointwise, dominated convergence gives (iv). For $h>0$, the difference quotient $[(X-K-h)^+-(X-K)^+]/h$ is bounded by $1$ in absolute value and converges to $-\mathbf 1\{X>K\}$ as $h\downarrow0$, so $c'_+(K)=-\P(X>K)$. In particular $c'_+(0)=-\P(X>0)\ge-1$.

\emph{Sufficiency.} On $(0,\infty)$ the convex function $c$ is continuous, its right derivative $c'_+$ is finite, nondecreasing and right-continuous, and $c(b)-c(a)=\int_a^b c'_+(u)\,du$ for $0<a<b$. At the end point $0$ convexity alone gives less: $c$ can jump down at $0$, with $c(0)>\lim_{K\downarrow0}c(K)$, and then $c'_+(0)=-\infty$. Condition (iii) excludes this jump. The chord slope $(c(h)-c(0))/h$ is nondecreasing in $h>0$ with limit $c'_+(0)\ge-1$, so $-h\le c(h)-c(0)\le h\,(c(1)-c(0))$ for $0<h\le1$. Hence $c$ is continuous at $0$ and Lipschitz on every $[0,b]$, so it is absolutely continuous there and the identity $c(b)-c(a)=\int_a^bc'_+(u)\,du$ also holds with $a=0$; and $c'_+$ is right-continuous at $0$, because $c'_+(0)\le c'_+(K)\le(c(K')-c(K))/(K'-K)$ for $0<K<K'$, and letting $K\downarrow0$, by continuity at $0$, and then $K'\downarrow0$ gives $c'_+(K)\to c'_+(0)$. If $c'_+(K_0)>0$ for some $K_0$, then $c(K)\ge c(K_0)+c'_+(K_0)(K-K_0)\to\infty$, contradicting (iv); so $c'_+\le0$, $c$ is nonincreasing and, by (iv), $c\ge0$. Let $L=\lim_{K\to\infty}c'_+(K)\le0$. Since $c'_+$ is nondecreasing, $c'_+\le L$ everywhere, so $c(K)\le F+LK$; $L<0$ would force $c<0$ for large $K$. Hence $L=0$.

Define $G(K)=0$ for $K<0$ and $G(K)=1+c'_+(K)$ for $K\ge0$. Then $G$ is nondecreasing and right-continuous, $G(0)=1+c'_+(0)\ge0$ by (iii), and $G(K)\to1$. It is therefore the distribution function of a random variable $X\ge0$ (Lebesgue--Stieltjes construction). By Tonelli's theorem,
\begin{align*}
\E(X-K)^+&=\int_K^\infty \P(X>u)\,du=\int_K^\infty\bigl(1-G(u)\bigr)du\\
&=-\int_K^\infty c'_+(u)\,du=c(K)-\lim_{b\to\infty}c(b)=c(K),
\end{align*}
and at $K=0$, where the last step uses the identity on $[0,b]$ and hence the continuity of $c$ at $0$ that (iii) provides, $\E X=c(0)=F$. Uniqueness holds because $c$ determines $c'_+$ and hence $G$.
\end{proof}

\begin{remark}
Extending $c$ by $c(K)=F-K$ for $K<0$, the second distributional derivative of $c$ on $\mathbb{R}$ is the law of $X$; the jump of $c'$ at $0$ carries the atom $\P(X=0)$. On a finite strike grid the necessary conditions reduce to slopes in $[-1,0]$ that are nondecreasing, which is the discrete test used in \texttt{src/qef/fx/arbitrage.py}. The implementation also applies the density condition on implied total variance of Gatheral and Jacquier (2014, eq.~(2.1) and Lemma~2.2).
\end{remark}

\section*{R3. Spanning}

The representation of a payoff by bonds, forwards and out-of-the-money options is the spanning formula of Carr and Madan (2001; first published 1998). The version below states explicit sufficient hypotheses, including the integrability condition under which expectations can be taken term by term. The proof is the standard Taylor expansion with integral remainder, written out under those hypotheses. In R3 and R4, $X$ is the terminal exchange rate and $\E$ and $\Q$ refer to the same pricing measure.

\begin{proposition}
Let $f:(0,\infty)\to\mathbb{R}$ have a locally absolutely continuous derivative $f'$, so that $f''$ exists almost everywhere and is locally integrable. Let $X>0$ almost surely with $\E X=F<\infty$, and write $P(K)=\E(K-X)^+$ and $C(K)=\E(X-K)^+$. If
\[
I=\int_0^F|f''(K)|\,P(K)\,dK+\int_F^\infty|f''(K)|\,C(K)\,dK<\infty,
\]
then $f(X)\in L^1$ and
\[
\E f(X)=f(F)+\int_0^F f''(K)\,P(K)\,dK+\int_F^\infty f''(K)\,C(K)\,dK.
\]
\end{proposition}

\begin{proof}
Fix $x>0$. Absolute continuity of $f$ and $f'$ on the interval between $x$ and $F$ gives $f(x)-f(F)=\int_F^x f'(u)\,du$ and $f'(u)=f'(F)+\int_F^u f''(K)\,dK$. Substituting and applying Fubini's theorem on the bounded interval (where $f''$ is integrable),
\[
f(x)=f(F)+f'(F)(x-F)+\int_F^x f''(K)(x-K)\,dK.
\]
For $x>F$ the last integral equals $\int_F^\infty f''(K)(x-K)^+dK$; for $x<F$ it equals $\int_x^F f''(K)(K-x)\,dK=\int_0^F f''(K)(K-x)^+dK$. Hence, for every $x>0$,
\[
f(x)=f(F)+f'(F)(x-F)+\int_0^F f''(K)(K-x)^+dK+\int_F^\infty f''(K)(x-K)^+dK.
\]
The map $(K,\omega)\mapsto f''(K)(K-X(\omega))^+$ is jointly measurable. By Tonelli's theorem, $\E\int_0^F|f''(K)|(K-X)^+dK=\int_0^F|f''(K)|P(K)\,dK$, and similarly for the call term, so both are at most $I<\infty$. Each term on the right is therefore integrable, hence so is $f(X)$, and Fubini's theorem allows expectation and integral to be exchanged. Since $\E(X-F)=0$, the formula follows.
\end{proof}

\begin{remark}
For $f(x)=\ln(x/F)$, $f''(K)=-K^{-2}$, which gives the log contract $\E\ln(X/F)=-\int_0^F P(K)K^{-2}dK-\int_F^\infty C(K)K^{-2}dK$. For $f(x)=(\ln(x/F))^2$, $f''(K)=2(1+\ln(F/K))/K^2$, which changes sign at $K=eF$; the absolute-value condition in the proposition covers such weights. The moment contracts of Bakshi, Kapadia and Madan (2003, Theorem 1) are of this type; they expand around the spot price with a constant interest rate, whereas the contracts here are centred at the forward $F$. Quotes observe $P$ and $C$ only between the $10\Delta$ strikes, so the tails of these integrals must be bounded (R4).
\end{remark}

\section*{R4. Bounds on the unobserved tails}

Jiang and Tian (2005, Proposition 2) bound the truncation error of model-free implied variance beyond the quoted strike range. The bounds below serve the same purpose under a different assumption, a power-type monotonicity of the distribution function in each tail, which makes them computable from the last quoted prices and extends them to the weights of the higher moment contracts. Lee (2004) relates the number of finite moments of the underlying to the limiting slope of implied variance in each wing, the asymptotic counterpart of the tail exponents used here. The lemma, the proposition and their proofs are my own; I have not found them in this form in the sources I consulted.

Let $G(u)=\Q(X\le u)$ and $\bar G(u)=\Q(X>u)$. By Tonelli's theorem, $P(K)=\int_0^K G(u)\,du$ and $C(K)=\int_K^\infty\bar G(u)\,du$. Let $K_{\min}<F<K_{\max}$ be the lowest and highest strikes at which prices are observed.

\begin{lemma}
(a) If $u\mapsto G(u)/u^\gamma$ is nondecreasing on $(0,K_{\min}]$ for some $\gamma>0$, then $P(K)\le P(K_{\min})(K/K_{\min})^{\gamma+1}$ for $0<K\le K_{\min}$. (b) If $u\mapsto u^\eta\bar G(u)$ is nonincreasing on $[K_{\max},\infty)$ for some $\eta>1$, then $C(K)\le C(K_{\max})(K/K_{\max})^{1-\eta}$ for $K\ge K_{\max}$.
\end{lemma}

\begin{proof}
(a) With $\lambda=K/K_{\min}\in(0,1]$ and the substitution $u=\lambda v$, $P(K)=\lambda\int_0^{K_{\min}}G(\lambda v)\,dv$. Since $\lambda v\le v$ and $G(w)/w^\gamma$ is nondecreasing, $G(\lambda v)\le\lambda^\gamma G(v)$. Hence $P(K)\le\lambda^{\gamma+1}\int_0^{K_{\min}}G(v)\,dv=\lambda^{\gamma+1}P(K_{\min})$. (b) With $\lambda=K/K_{\max}\ge1$ and $u=\lambda v$, $C(K)=\lambda\int_{K_{\max}}^\infty\bar G(\lambda v)\,dv$, and $\bar G(\lambda v)\le\lambda^{-\eta}\bar G(v)$ because $w^\eta\bar G(w)$ is nonincreasing. Hence $C(K)\le\lambda^{1-\eta}C(K_{\max})$.
\end{proof}

\begin{proposition}
Under the hypotheses of the lemma, for any measurable $w\ge0$ on $(0,K_{\min}]$,
\[
\int_0^{K_{\min}}w(K)P(K)\,dK\le\frac{P(K_{\min})}{K_{\min}^{\gamma+1}}\int_0^{K_{\min}}w(K)K^{\gamma+1}dK.
\]
In particular $\int_0^{K_{\min}}P(K)K^{-2}dK\le P(K_{\min})/(\gamma K_{\min})$. For $w(K)=2(1+\ln(F/K))K^{-2}$,
\[
\int_0^{K_{\min}}w(K)P(K)\,dK\le\frac{2P(K_{\min})}{K_{\min}}\Bigl[\frac{1+\ln(F/K_{\min})}{\gamma}+\frac{1}{\gamma^2}\Bigr].
\]
For the upper tail and any measurable $w$ on $[K_{\max},\infty)$,
\[
\int_{K_{\max}}^\infty|w(K)|\,C(K)\,dK\le C(K_{\max})K_{\max}^{\eta-1}\int_{K_{\max}}^\infty|w(K)|\,K^{1-\eta}dK;
\]
in particular $\int_{K_{\max}}^\infty C(K)K^{-2}dK\le C(K_{\max})/(\eta K_{\max})$, and for $w(K)=2(1-\ln(K/F))K^{-2}$ with $F<K_{\max}\le eF$ and $\ell=\ln(K_{\max}/F)$,
\[
\int_{K_{\max}}^\infty|w(K)|\,C(K)\,dK\le\frac{2C(K_{\max})}{K_{\max}}\Bigl[\frac{1-\ell}{\eta}-\frac{1}{\eta^2}+\frac{2e^{-\eta(1-\ell)}}{\eta^2}\Bigr].
\]
\end{proposition}

\begin{proof}
The first inequality is the lemma multiplied by $w\ge0$ and integrated. For $w(K)=K^{-2}$, $\int_0^{K_{\min}}K^{\gamma-1}dK=K_{\min}^\gamma/\gamma$. For the logarithmic weight, integration by parts gives $\int_0^{K_{\min}}K^{\gamma-1}\ln(F/K)\,dK=K_{\min}^\gamma\ln(F/K_{\min})/\gamma+K_{\min}^\gamma/\gamma^2$, the boundary term at zero vanishing because $K^\gamma\ln(F/K)\to0$. The upper-tail bounds follow from part (b) of the lemma. For $w(K)=K^{-2}$, $\int_{K_{\max}}^\infty K^{-1-\eta}dK=K_{\max}^{-\eta}/\eta$. For the logarithmic weight, the substitution $u=\ln(K/F)$ gives $F^{-\eta}\int_\ell^\infty|1-u|e^{-\eta u}du$, which is split at $u=1$ and integrated by parts.
\end{proof}

\begin{remark}
The bounds use only observed wing prices and an assumed tail exponent. A tail term whose weight has constant sign lies between $0$ and its bound, taken with the sign of the weight. A sign-changing weight is split into its positive and negative parts, $w=w_+-w_-$ with $w_\pm=\max(\pm w,0)$; by the proposition the tail integrals against $w_+$ and $w_-$ lie in $[0,B_+]$ and $[0,B_-]$, where $B_\pm$ is the bound for $w_\pm$, so the term lies in $[-B_-,B_+]$. The interval $[-B,B]$ with $B=B_++B_-$, the bound for $|w|$, is also valid but coarser. Because $\gamma$ and $\eta$ are not observed, risk-neutral variance and skewness are reported as intervals over a stated range of exponents, obtained by propagating the intervals of the individual contracts, never as point values. Between $K_{\min}$ and $K_{\max}$, $P$ and $C$ are taken from the calibrated smile and the integrals are computed by composite Simpson's rule, with $F$ a node and each side integrated separately, because the out-of-the-money integrand has a kink at $F$ where $P'-C'=1$. With node spacing $h$ on a side $[a,b]$ and an even number of subintervals for Simpson's rule, the composite trapezium error is at most $(b-a)h^2\max|g''|/12$ if $g\in C^2[a,b]$, and the composite Simpson error at most $(b-a)h^4\max|g^{(4)}|/180$ if $g\in C^4[a,b]$ (S\"uli and Mayers, 2003).
\end{remark}

\section*{R5(a). Strike from pips delta at a flat volatility}

This result is standard. The pips spot and forward deltas are those defined by Reiswich and Wystup (2012; working-paper version of 20 March 2010, Sections 1.3 and 1.5), and their inversion at a single volatility is elementary; the proof is written out for completeness, and none of it is my own. It is the closed form used for pips deltas in \texttt{src/qef/fx/gk.py} and for the strikes $K^{ms}_{c,p}$ in the proof of R8.

\begin{proposition}
Let $\sigma>0$, let $D>0$ be the base-currency discount factor (with forward delta, $D=1$), and let $\phi=1$ for a call and $\phi=-1$ for a put. The pips delta at the flat volatility $\sigma$, $\Delta(K)=\phi D\,\Phi\bigl(\phi\,d_+(K)\bigr)$, with $s=\sigma\sqrt\tau$ in $d_+$, is continuous and strictly decreasing in $K$, and maps $(0,\infty)$ onto $(0,D)$ for a call and onto $(-D,0)$ for a put. For each $\Delta$ with the sign of $\phi$ and $0<|\Delta|<D$, the unique strike with that delta is
\[
K=F\exp\Bigl(\tfrac12\sigma^2\tau-\phi\,\sigma\sqrt\tau\,\Phi^{-1}\bigl(|\Delta|/D\bigr)\Bigr).
\]
\end{proposition}

\begin{proof}
Since $\Phi$ takes values in $(0,1)$, $\Delta(K)$ has the sign of $\phi$ and $|\Delta(K)|=D\,\Phi(\phi\,d_+(K))\in(0,D)$. The map $K\mapsto d_+(K)$ is a strictly decreasing bijection of $(0,\infty)$ onto $\mathbb R$, and $\Delta'(K)=\phi^2D\,\varphi(\phi\,d_+)\,d_+'(K)=D\,\varphi(d_+)\,d_+'(K)<0$, since $\varphi$ is even. As $K\downarrow0$, $d_+\to\infty$, so $\Delta\to D$ for a call and $\Delta\to0$ for a put; as $K\to\infty$, $d_+\to-\infty$, so $\Delta\to0$ for a call and $\Delta\to-D$ for a put. Continuity and strict monotonicity give the stated ranges and uniqueness. Given $\Delta$ with the sign of $\phi$ and $0<|\Delta|<D$, the equation $\Phi(\phi\,d_+)=|\Delta|/D$ gives $\phi\,d_+=\Phi^{-1}(|\Delta|/D)$, hence $d_+=\phi\,\Phi^{-1}(|\Delta|/D)$ because $\phi^2=1$. Substituting into $\ln(F/K)=s\,d_+-s^2/2$ gives $\ln(K/F)=s^2/2-\phi\,s\,\Phi^{-1}(|\Delta|/D)$, which is the stated formula.
\end{proof}

\section*{R5(b). Premium-adjusted call delta}

The premium-adjusted delta, its non-monotonicity in the strike and the market practice of taking the strike to the right of the delta maximum are described by Reiswich and Wystup (2012; working-paper version of 20 March 2010, Sections 1.3 and 1.5), who also show that the premium-adjusted delta lies below the simple delta because the call price is positive, and state that the maximum solves $\sigma\sqrt\tau\,\Phi(d_-)=\varphi(d_-)$. J\"ackel (2020, Section 2, eq.~(20)) shows that the premium-adjusted forward delta of a call has a single maximum, located where the inverse Mills ratio equals the total volatility, and takes the larger of the two solutions, which is the strike on the conventional branch. The proposition below states this shape in the notation used here; its proof is the same inverse Mills ratio argument, written out for completeness, and none of it is my own.

Let $D>0$ and $h(K)=D\,(K/F)\,\Phi\bigl(d_-(K)\bigr)$ for $K>0$. With spot delta, $D$ is the base-currency discount factor; with forward delta, $D=1$.

\begin{lemma}
The inverse Mills ratio $\lambda(d)=\varphi(d)/\Phi(d)$ is strictly decreasing on $\mathbb{R}$, with $\lambda(d)\to\infty$ as $d\to-\infty$ and $\lambda(d)\to0$ as $d\to\infty$.
\end{lemma}

\begin{proof}
Differentiation gives $\lambda'(d)=-\lambda(d)\,[d+\lambda(d)]$, so it suffices that $d\,\Phi(d)+\varphi(d)>0$. This function has derivative $\Phi(d)>0$ and tends to $0$ as $d\to-\infty$; hence $d\,\Phi(d)+\varphi(d)=\int_{-\infty}^d\Phi(u)\,du>0$. The limits follow from $\Phi\to0$, $\varphi\to0$ and the Mills-ratio asymptotics $\Phi(d)\sim\varphi(d)/|d|$ as $d\to-\infty$. The lemma is known: the normal density is log-concave, so $\Phi$ is log-concave (Bagnoli and Bergstrom, 2005, Section 2.1, Theorem 1) and $\lambda=(\log\Phi)'$ is decreasing, and inequalities of this type for the Mills ratio are classical (Sampford, 1953). The short direct proof above is written out for completeness.
\end{proof}

\begin{proposition}
The function $h$ is continuous on $(0,\infty)$, tends to $0$ as $K\downarrow0$ and as $K\to\infty$, and satisfies $h<D$. There is a unique $K^*$ such that $h$ is strictly increasing on $(0,K^*]$ and strictly decreasing on $[K^*,\infty)$; $K^*=F\exp(-d^*s-s^2/2)$, where $d^*$ is the unique solution of $\lambda(d)=s$. Each $\Delta\in(0,h(K^*))$ has exactly two preimages, one on each side of $K^*$.
\end{proposition}

\begin{proof}
As $K\downarrow0$, $d_-\to\infty$, so $h\to0$. As $K\to\infty$, $d_-\to-\infty$ and, by the Mills-ratio asymptotics and the identity above, $K\Phi(d_-)\sim K\varphi(d_-)/|d_-|=F\varphi(d_+)/|d_-|\to0$. Moreover $F\Phi(d_+)-K\Phi(d_-)$ is the undiscounted call price, which is positive, so $h<D\,\Phi(d_+)<D$.

Since $d_-'(K)=-1/(Ks)$,
\[
h'(K)=\frac{D}{F}\Bigl[\Phi(d_-)-\frac{\varphi(d_-)}{s}\Bigr]=\frac{D\,\Phi(d_-)}{F\,s}\,\bigl[s-\lambda(d_-)\bigr].
\]
By the lemma, $s-\lambda(d)$ is strictly increasing in $d$ and changes sign exactly once, at $d^*$. Because $d_-$ is strictly decreasing in $K$, $h'>0$ for $K<K^*$ and $h'<0$ for $K>K^*$, where $d_-(K^*)=d^*$. The stated formula for $K^*$ inverts $d_-$. The final claim follows from the intermediate value theorem and strict monotonicity on each side.
\end{proof}

\begin{remark}
The market convention selects the preimage in $[K^*,\infty)$. At one month and $25\Delta$ this root lies far to the right of $K^*$; bracketing matters at long maturities and high volatilities.
\end{remark}

\section*{R5(c). Delta-to-strike inversion with a smile}

Reiswich (2010, Section 3.8 and Fig.~3.6) shows that the delta implied by a smile need not be monotone in the strike, so that a quoted delta can correspond to several strikes or to none. Fukasawa (2012, Lemma 2.1 and Propositions 2.1--2.2) shows, by the same derivative calculation as below, that $d_+$ and $d_-$ are strictly decreasing in the strike for every arbitrage-free smile, and Mingone (2023, Lemma 4.1) links this monotonicity to a unique strike for each delta. The part that is my own is only the explicit slope bound in the proposition, which can be checked on a fitted smile that is not assumed to be free of arbitrage, stated here for pips spot and forward deltas with the discount factor $D$.

\begin{proposition}
Let $\sigma:(0,\infty)\to(0,\infty)$ be continuously differentiable and let $\Delta_\phi(K)=\phi D\,\Phi\bigl(\phi\,d_+(K,\sigma(K))\bigr)$ be the pips spot delta ($D$ the base-currency discount factor) or forward delta ($D=1$) of a call ($\phi=1$) or put ($\phi=-1$). On an interval $J$, if
\[
|\sigma'(K)|\,K\sqrt\tau\,|d_-(K,\sigma(K))|<1\qquad\text{for all }K\in J,
\]
then $\Delta_\phi$ is strictly decreasing on $J$, and each delta attained on $J$ is attained at exactly one strike in $J$.
\end{proposition}

\begin{proof}
With $d_+=\ln(F/K)/(\sigma\sqrt\tau)+\sigma\sqrt\tau/2$, one has $\partial_Kd_+=-1/(K\sigma\sqrt\tau)$ and $\partial_\sigma d_+=-d_-/\sigma$. By the chain rule,
\[
\frac{d}{dK}\,d_+\bigl(K,\sigma(K)\bigr)=-\frac{1}{K\sigma\sqrt\tau}\Bigl[1+K\sqrt\tau\,d_-\,\sigma'(K)\Bigr],
\]
which is negative on $J$ under the stated condition. Since $\Delta_\phi'(K)=\phi^2D\,\varphi(d_+)\,\frac{d}{dK}d_+=D\,\varphi(d_+)\,\frac{d}{dK}d_+$, the function $\Delta_\phi$ is strictly decreasing on $J$, and injectivity gives uniqueness.
\end{proof}

\begin{remark}
In log-moneyness $k=\ln(K/F)$, $K\sigma'(K)=\partial\sigma/\partial k$, so the condition reads $|\partial\sigma/\partial k|\sqrt\tau\,|d_-|<1$. Across the calibrated one-month smiles of the primary sample, its maximum over $\pm4$ ATM standard deviations is $0.76$ (\texttt{scripts/check\_reported\_conditions.py}), so the smile delta-to-strike map is injective on that range at every month-end. For premium-adjusted delta $\phi D(K/F)\Phi(\phi d_-)$ the analogous derivative is $\phi(D/F)\bigl[\Phi(\phi d_-)+\phi K\varphi(d_-)\,\tfrac{d}{dK}d_-\bigr]$ with $\tfrac{d}{dK}d_-=-\bigl[1+K\sqrt\tau\,d_+\sigma'(K)\bigr]/(K\sigma\sqrt\tau)$; numerically, the premium-adjusted put delta is strictly monotone and the call delta strictly decreasing to the right of its maximum on the same range at every month-end.
\end{remark}

\section*{R6. Diffusive null of the hedge cost}

The comparison of hedged and unhedged carry returns follows Jurek (2014), who reads the resulting ratio as an upper bound on the crash share and notes that an unlevered option hedge also gives up part of the diffusive premium (working-paper version of August 2013, Section 2.2). Farhi et al. (2015, Section 5.1, eqs. (8)--(9)) state the leading-order version of the result below: a put of delta $\Delta_P$ leaves a fraction $1+\Delta_P$ of the Gaussian risk premium. The mean-value form and the error bound below are my own short derivation.

Let the undiscounted put value with lognormal forward $G$ and volatility $\sigma$ be
\[
p(G)=K\,\Phi\bigl(-d_-(G)\bigr)-G\,\Phi\bigl(-d_+(G)\bigr),\qquad d_\pm(G)=\frac{\ln(G/K)\pm s^2/2}{s}.
\]
Using $K\varphi(d_-)=G\varphi(d_+)$, $p'(G)=-\Phi(-d_+(G))$.

Suppose that under $\Q$, $S_\tau$ is lognormal with mean $F$ and log-variance $s^2$, and that under $\P$ it is lognormal with the same log-variance and mean $Fe^{\lambda\tau}$ (no jumps, no volatility premium). A unit long forward costs nothing and has $\P$-expected payoff $\mu_U=F(e^{\lambda\tau}-1)$. Adding a put bought at its $\Q$-value, with the premium financed to $\tau$, gives $\mu_H=\mu_U+p(Fe^{\lambda\tau})-p(F)$.

\begin{proposition}
If $\lambda\ne0$, the hedge-cost ratio $\theta_0=1-\mu_H/\mu_U$ satisfies $\theta_0=\Phi(-d_+(\xi))$ for some $\xi$ between $F$ and $Fe^{\lambda\tau}$, and
\[
\bigl|\theta_0-\Phi(-d_1)\bigr|\le\varphi(0)\,\frac{|\lambda|\sqrt{\tau}}{\sigma},\qquad d_1=d_+(F).
\]
\end{proposition}

\begin{proof}
By the mean value theorem, $p(Fe^{\lambda\tau})-p(F)=-\Phi(-d_+(\xi))\,F(e^{\lambda\tau}-1)$ for some $\xi$ between the two forwards, so $\theta_0=-[p(Fe^{\lambda\tau})-p(F)]/\mu_U=\Phi(-d_+(\xi))$. Since $|d_+(\xi)-d_+(F)|=|\ln(\xi/F)|/s\le|\lambda|\tau/s$ and $\Phi$ is Lipschitz with constant $\varphi(0)$, the bound follows.
\end{proof}

\begin{remark}
With no crash premium at all, hedging with a put of forward delta $-\Phi(-d_1)$ surrenders approximately that fraction of the carry premium: about $0.25$ for a $25\Delta$ put and $0.10$ for a $10\Delta$ put. This is the null against which the realised ratio $\theta_{UB}$ and the skew component of E3 are read.
\end{remark}

\section*{R7. Entropy decomposition of currency premia}

The decomposition is due to Backus, Foresi and Telmer (2001): with the exchange rate equal to the ratio of the two pricing kernels (their Proposition 1), the expected excess return on a forward position is the difference between the log of the expectation and the expectation of the log of the two kernels (their eq.~(11)), and its higher cumulants enter only through that difference (eqs.~(12)--(13)). Conditional entropy, as this difference is now called, and the entropy bound are as in Backus, Chernov and Zin (2014), who relate the bound to Bansal and Lehmann (1997) and Alvarez and Jermann (2005). Lustig and Verdelhan (2019, Section 2.5) restate the decomposition with a wedge between the exchange rate and the ratio of kernels, and Jiang, Krishnamurthy and Lustig (2021) link the value of the dollar to a convenience yield inferred from the Treasury basis, a covered-parity deviation measured with government bond yields. The identity below writes the same decomposition with the covered-interest-parity basis as the wedge on the forward side; the derivation is elementary.

Fix $t$ and write $\E_t$ for $\E[\cdot\mid\mathcal F_t]$. For a positive random variable $M$ with $\E_t|\log M|<\infty$, the conditional entropy is $L_t(M)=\log\E_tM-\E_t\log M\ge0$. Let $M,M^*>0$ be $\mathcal F_{t+1}$-measurable discount factors in $L^2$ that price domestic and foreign payoffs in their own currencies, with $\E_tM=e^{-r_t}$, $\E_tM^*=e^{-r^*_t}$, $\E_t|\log M|,\E_t|\log M^*|<\infty$ and $MS_{t+1}/S_t\in L^2$, where $S$ is the domestic price of foreign currency. Write $m=\log M$, $\Delta s=\log(S_{t+1}/S_t)$, and $x_t=(f_t-s_t)-(r_t-r^*_t)$ for the covered-interest-parity basis. With exchange rates quoted in foreign currency per USD, as in Du, Tepper and Verdelhan (2018; NBER Working Paper 23170, eq. (4)), $x_t$ is minus their basis.

\begin{lemma}
Suppose every foreign payoff $X^*\in L^2(\mathcal F_{t+1})$ is traded and priced consistently in both currencies, $S_t\E_t[M^*X^*]=\E_t[MS_{t+1}X^*]$. Then $M^*=MS_{t+1}/S_t$ almost surely.
\end{lemma}

\begin{proof}
The difference $D=M^*-MS_{t+1}/S_t$ lies in $L^2$ and satisfies $\E_t[DX^*]=0$ for all $X^*\in L^2$. Taking $X^*=D$ gives $\E_tD^2=0$.
\end{proof}

\begin{proposition}
Under the lemma's conclusion, (i) $\E_t[\Delta s+r^*_t-r_t]=L_t(M)-L_t(M^*)$ and $\E_t[s_{t+1}-f_t]=L_t(M)-L_t(M^*)-x_t$. (ii) If $k(u)=\log\E_te^{um}$ is finite near $0$ and its Maclaurin series has radius of convergence greater than $1$, then $L_t(M)=\sum_{j\ge2}\kappa_{j,t}/j!$, where $\kappa_{j,t}$ are the conditional cumulants of $m$. (iii) For any gross return $R>0$ with $\E_t[MR]=1$ and $\E_t|\log R|<\infty$, $L_t(M)\ge\E_t\log R-r_t$.
\end{proposition}

\begin{proof}
(i) By definition, $\E_tm=-r_t-L_t(M)$ and $\E_tm^*=-r^*_t-L_t(M^*)$. The lemma gives $\Delta s=m^*-m$, so $\E_t\Delta s=r_t-r^*_t+L_t(M)-L_t(M^*)$. The second identity follows from $s_{t+1}-f_t=\Delta s-(f_t-s_t)=\Delta s-(r_t-r^*_t)-x_t$. (ii) $L_t(M)=k(1)-k'(0)$ with $k'(0)=\kappa_{1,t}$, and $k(1)=\sum_{j\ge1}\kappa_{j,t}/j!$ under the stated convergence. (iii) By Jensen's inequality, $0=\log\E_t[MR]\ge\E_tm+\E_t\log R$, so $L_t(M)=-r_t-\E_tm\ge\E_t\log R-r_t$.
\end{proof}

\begin{remark}
The identity in (i) needs only that some pair of discount factors prices bonds and exchange rates consistently; completeness makes the pair unique. Brandt, Cochrane and Santa-Clara (2006) show that the smoothness of exchange rates relative to the volatility of marginal utility implied by equity premia requires a very high degree of international risk sharing. Options on $S$ reveal the risk-neutral law of $\Delta s=m^*-m$, not the separate cumulants of $m$ and $m^*$, so they constrain but do not identify the split of $L_t(M)-L_t(M^*)$ across cumulant orders. Hedged and unhedged carry returns give two lower bounds on the same $L_t(M)$ through (iii).
\end{remark}

\section*{R8. Local well-posedness of the smile calibration}

Reiswich (2010, Theorems 2 and 3 and Appendices A and B) uses the implicit function theorem to show that a volatility-strike function exists near the at-the-money point and that the market-strangle calibration of the simplified parabolic smile is locally solvable for the smile strangle near a zero risk reversal. The proposition below applies the same argument to the three SABR parameters and the three quotes. Le Floc'h and Kennedy (2014, Section 4, eqs.~(30)--(31)) invert the leading-order expansion of the Hagan formula (Hagan et al., 2002, eq.~(3.1a)) in closed form, recovering $(\alpha,\rho,\nu)$ from the at-the-money level, skew and curvature, which is an approximate, global counterpart of the local and exact statement here. The three calibration equations are the conditions of Reiswich and Wystup (2012; working-paper version, Section 2, eqs.~(17), (18) and (25)), to which Reiswich (2010, Section 3.2) calibrates SABR by least squares, and Henrard (2011) uses the implicit function theorem in the same way for any calibration with as many instruments as parameters. In the lemma, the value of $z/x(z)$ at $z=0$ and its expansion are known (Choi and Wu, 2021, eqs.~(4)--(5)), and the integral representation is Hadamard's lemma (Milnor, 1963, Section 2, Lemma 2.1). What I claim as my own is the explicit statement that the $\beta=1$ formula is jointly smooth on the stated domain, including $\nu=0$, and the local well-posedness statement for SABR under conditions (i) to (iv), written for FX delta conventions; I have not found these in this form in the sources I consulted.

\begin{lemma}
For $\beta=1$ the Hagan, Kumar, Lesniewski and Woodward (2002) implied volatility $\sigma(K;\theta)$, $\theta=(\alpha,\rho,\nu)$, is infinitely differentiable on $\{K>0\}\times(0,\infty)\times(-1,1)\times[0,\infty)$.
\end{lemma}

\begin{proof}
For $\beta=1$, $\sigma=\alpha\,\bigl(z/x(z)\bigr)\bigl[1+\bigl(\rho\nu\alpha/4+(2-3\rho^2)\nu^2/24\bigr)\tau\bigr]$ with $z=(\nu/\alpha)\ln(F/K)$ and $x(z)=\ln\bigl[(\sqrt{1-2\rho z+z^2}+z-\rho)/(1-\rho)\bigr]$. Since $(1-2\rho z+z^2)-(z-\rho)^2=1-\rho^2>0$, the square root exceeds $|z-\rho|$, so the argument of the logarithm is positive and $x$ is smooth. Moreover $x'(z)=(1-2\rho z+z^2)^{-1/2}>0$ and $x(0)=0$, so $x(z,\rho)=z\,h(z,\rho)$ with $h(z,\rho)=\int_0^1(1-2\rho tz+t^2z^2)^{-1/2}\,dt$. Since $1-2\rho tz+t^2z^2\ge1-\rho^2>0$, the integrand is smooth and positive on $\mathbb R\times(-1,1)\times[0,1]$, so by differentiation under the integral $h$ is smooth and positive, and $z/x(z)=1/h(z,\rho)$ (equal to $1$ at $z=0$) is jointly smooth in $(z,\rho)$. Compositions and products of smooth maps are smooth.
\end{proof}

Fix $F,\tau,D$ and the delta convention, and write $q=(\sigma_{ATM},RR,BF)$ for the quotes. For parameters $\theta$, let $K_A(\theta)$ solve $\ln(K/F)\mp\sigma(K;\theta)^2\tau/2=0$ (the delta-neutral-straddle strike, with the sign given by the premium convention). Let $K_c(\theta)$ and $K_p(\theta)$ solve $\Delta\bigl(K,\sigma(K;\theta)\bigr)=\pm0.25$. Let $K^{ms}_{c}(q)$ and $K^{ms}_{p}(q)$ be the $\pm25\Delta$ strikes at the single volatility $\sigma_{ATM}+BF$. Define $\Psi=(\Psi_1,\Psi_2,\Psi_3)$ by
\[
\Psi_1=\sigma(K_A;\theta)-\sigma_{ATM},\qquad \Psi_2=\sigma(K_c;\theta)-\sigma(K_p;\theta)-RR,
\]
\[
\Psi_3=\sum_{i\in\{c,p\}}\bigl[V\bigl(K^{ms}_i,\sigma(K^{ms}_i;\theta)\bigr)-V\bigl(K^{ms}_i,\sigma_{ATM}+BF\bigr)\bigr],
\]
where $V$ is the undiscounted Garman--Kohlhagen premium of the corresponding option. The calibrated smile solves $\Psi(\theta,q)=0$. The market-strangle condition is that of Reiswich and Wystup (2012). For $\theta$ near $\theta^*$, $K_A(\theta)$, $K_c(\theta)$ and $K_p(\theta)$ denote the $C^1$ branches through the strikes at $\theta^*$ given by the implicit function theorem under (i) and (ii) below; for a premium-adjusted call, the strike at $\theta^*$ is the root with $\partial_K\Delta<0$ (the conventional branch), which is the root that \texttt{strike\_from\_delta\_smile} selects.

\begin{proposition}
Let $\Psi(\theta^*,q^*)=0$. Suppose that at $(\theta^*,q^*)$: (i) $1\mp K\sigma\sigma'(K)\tau\ne0$ at $K_A$; (ii) the strike derivative of $\Delta(K,\sigma(K;\theta^*))$ is non-zero at $K_c$ and $K_p$ (for pips delta, for instance under the condition of R5(c); for premium-adjusted delta the derivative in the remark after R5(c) must be checked directly, since a flat smile satisfies the R5(c) condition while the call derivative vanishes at $K^*$); (iii) for premium-adjusted delta, the $25\Delta$ call strike at the flat volatility $\sigma^*_{ATM}+BF^*$ lies strictly to the right of the maximiser, $K^{ms}_c>K^*$ (equivalently $0.25<h(K^*)$ in the notation of R5(b)); and (iv) $D_\theta \Psi(\theta^*,q^*)$ is non-singular. Then there are neighbourhoods $U\ni q^*$ and $W\ni\theta^*$ and a unique $C^1$ map $\theta:U\to W$ with $\Psi(\theta(q),q)=0$. In particular $\theta^*$ is the only calibration in $W$ with strikes on these branches, and the calibrated smile varies smoothly with the quotes.
\end{proposition}

\begin{proof}
Since $x$ has the sign of $z$, $z/x(z)>0$, so $\sigma(\cdot;\theta)>0$ for all $K$ exactly when $1+\bigl(\rho\nu\alpha/4+(2-3\rho^2)\nu^2/24\bigr)\tau>0$. This holds at $\theta^*$ because $\sigma(K_A;\theta^*)=\sigma_{ATM}>0$, hence near $\theta^*$, and $\Delta$ and $V$ are evaluated only at positive volatilities. By the lemma, each defining equation of $K_A$, $K_c$, $K_p$ is then $C^1$ in $(K,\theta)$. Condition (i) is the non-vanishing of the $K$-derivative of the first; condition (ii) that of the other two. The implicit function theorem therefore makes $K_A$, $K_c$ and $K_p$ $C^1$ functions of $\theta$ near $\theta^*$. The strikes $K^{ms}_i$ depend on $q$ through $\sigma_{ATM}+BF$: in closed form for pips delta, $K^{ms}_{c,p}=F\exp\bigl(\tfrac12s_q^2\mp s_q\Phi^{-1}(0.25/D)\bigr)$ with $s_q=(\sigma_{ATM}+BF)\sqrt\tau$, by R5(a); and by the implicit function theorem for premium-adjusted delta, since $h'<0$ on $(K^*,\infty)$ by the proof of R5(b) and condition (iii), and $\partial_K\bigl[-D(K/F)\Phi(-d_-)\bigr]=-(D/F)\bigl[\Phi(-d_-)+\varphi(d_-)/s\bigr]<0$ for the put. Hence $\Psi$ is $C^1$ in $(\theta,q)$ near $(\theta^*,q^*)$, and the implicit function theorem applied to $\Psi$ with (iv) gives the claim.
\end{proof}

\begin{remark}
Condition (iv) is checked at every calibration through the finite-difference Jacobian $J$ returned by the solver. $J$ is taken in $x=(\ln\alpha,\operatorname{atanh}\rho,\ln\nu)$, with $\Psi_3$ divided by the flat-volatility strangle vega $\mathcal V$, so $J=\operatorname{diag}(1,1,\mathcal V^{-1})\,D_\theta\Psi\,\operatorname{diag}(\alpha,1-\rho^2,\nu)$, which is non-singular exactly when $D_\theta\Psi$ is. Across the primary sample $\operatorname{cond}(J)$ is at most 125 under the market-strangle reading and at most 634 under the smile-strangle reading (\texttt{scripts/check\_reported\_conditions.py}). Under the smile-strangle reading, $\Psi_3$ is replaced by $\tfrac12[\sigma(K_c;\theta)+\sigma(K_p;\theta)]-\sigma_{ATM}-BF$, and the proposition and its proof hold without condition (iii). The proposition is local: it does not assert existence of a solution for arbitrary quotes, and it does not exclude other solutions far from $\theta^*$. The implementation therefore restarts from several points and records non-convergence.
\end{remark}

\section*{R10. Admissible tail exponents and sharp identification of moments}

R4 bounds each moment contract separately, and it leaves open which exponents are compatible with the quotes. This section answers both questions under the hypotheses of R4. The notation is that of R3 and R4, with $X$ the terminal exchange rate under the quote-currency forward measure, $P_0=P(K_{\min})$, $G_0=G(K_{\min})=P'(K_{\min}+)$, $C_0=C(K_{\max})$ and $\bar G_0=\Q(X\ge K_{\max})=-C'(K_{\max}-)$, all taken from the calibrated smile on $[K_{\min},K_{\max}]$. The one-sided slopes are those inside the observed range, which the smile determines: by the difference-quotient argument in the proof of R2, $P'(K_{\min}+)=\Q(X\le K_{\min})$ and $C'(K_{\max}-)=-\Q(X\ge K_{\max})$. An atom of $X$ at either boundary strike is therefore not excluded by the observations, and $\bar G_0=\bar G(K_{\max})$ when $X$ has no atom at $K_{\max}$. Write $H_L(\gamma)$ for the hypothesis that $G(u)/u^\gamma$ is nondecreasing on $(0,K_{\min}]$ and $H_U(\eta)$, $\eta>1$, for the hypothesis that $u^\eta\bar G(u)$ is nonincreasing on $[K_{\max},\infty)$.

Several ingredients are known, and I credit them where they are used. $H_L(\gamma)$ is the restriction to the lower tail of the reversed hazard rate order, equivalently the star order, of $X$ against the power-function law, and $H_U(\eta)$ is the corresponding statement against the Pareto law (Lando, Arab and Oliveira, 2022, Definition 2); with a density, $H_L(\gamma)$ says that the elasticity $ug(u)/G(u)$ of the distribution function, studied by Veres-Ferrer and Pavia (2022, Sections 2--5), is at least $\gamma$ in the tail. The quantity $K_{\min}P'(K_{\min})/P(K_{\min})$ is proposed as the natural exponent of a power-law put extrapolation by Benaim, Dodgson and Kainth (2008, footnote 6 of the version consulted), and Qin, Yang, Che and Feng (2026, Propositions 1 and 2) construct the pure power-law tails that match the boundary price and slope. Taleb, Yarckin, Mann, Delic and Spitznagel (2023, Section IV) observe that convexity at the strike where a Pareto tail is attached restricts its index, which is close to the admissibility statement in (a) below, and Bollerslev, Todorov and Xu (2015, eq.~(3.3)) estimate tail decay from log-log slopes of out-of-the-money prices. The reduction of bounds over a class generated by a family of distributions to a moment problem over the mixing measure follows Popescu (2005, Section 3); that extreme points of the resulting set are mixtures of at most two generators is Theorem 2.1 of Winkler (1988), as quoted by Pinelis (2012, Theorem 12); and the description of the bounds as sections of the convex hull of a curve is that of Kemperman (1968, Section 0 and Theorem 1). The breakdown point in (e) below is the identification breakdown point of Horowitz and Manski (1995), generalised by Stoye (2005, 2010), as set out by Masten and Poirier (2020, pp.~42 and~53), and it is defined, as in Kline and Santos (2013, Example 2.3), as the level of an assumption at which the sign of a parameter stops being identified. Zhuang (2026) studies sharp bounds on option-implied tail quantities from finitely many quotes without a tail-shape hypothesis, and notes that the boundary slope $G_0$ is itself only set-identified from finite quotes; here it is taken from the calibrated smile.

What I claim as my own is the following. In (a), the observation that convexity at the strike where a Pareto tail is attached restricts its index is that of Taleb et al.\ (2023), it is close to the admissibility statement, and it is not mine; what I claim is the characterisation of the exponents compatible with the observed boundary price and slope under $H_L$ and $H_U$, by a condition that is necessary and sufficient within the whole class, together with the uniqueness of the tail within the whole class at the largest exponent. I also claim the generating family of the class with the boundary value fixed, the resulting description of the identified set of the moment contracts, and its consequences for R4; the check of setting (i) in (d); and the use of the breakdown point for the sign of option-implied skewness.

\begin{lemma}
Let $\gamma>0$. A nondecreasing right-continuous function $G$ on $(0,K_{\min}]$ with $G(K_{\min})=G_0$ satisfies $H_L(\gamma)$ if and only if there is a probability measure $\mu$ on $[0,1]$ with
\[
G(K_{\min}z)=G_0\,z^\gamma\,\mu([0,z]),\qquad 0<z<1.
\]
Then, for measurable $\phi$ with $\int_0^{K_{\min}}G|\phi|\,du<\infty$,
\[
\int_0^{K_{\min}}G(u)\phi(u)\,du=G_0K_{\min}\int_{[0,1]}\int_{z_c}^1z^\gamma\phi(K_{\min}z)\,dz\,\mu(dz_c).
\]
\end{lemma}

\begin{proof}
Under $H_L(\gamma)$, $\bar h(z)=G(K_{\min}z)/(G_0z^\gamma)$ is nondecreasing and right-continuous on $(0,1]$ with $0\le\bar h\le\bar h(1)=1$. Let $\mu$ be the measure with $\mu([0,z])=\bar h(z)$ for $0\le z<1$, with $\mu(\{0\})=\bar h(0+)$, and with the remaining mass $1-\bar h(1-)$ at $1$. Conversely, for any such $\mu$ the function defined by the display is nondecreasing, right-continuous, at most $G_0$ and satisfies $H_L(\gamma)$, because it is a product of the nonnegative nondecreasing functions $z^\gamma$ and $\mu([0,z])$; it tends to $0$ at $0$, so the law it describes has no mass at $0$. The integral identity is Tonelli's theorem applied to $\mu([0,z])=\int\mathbf 1\{z_c\le z\}\,\mu(dz_c)$.
\end{proof}

The upper tail has the same structure, with $K_{\max}$, $\bar G_0$, the power $z^{-\eta}$ and a probability measure $\nu$ on $[1,\infty]$ in place of $K_{\min}$, $G_0$, $z^\gamma$ and $\mu$. Let $\eta>1$ and $\bar G_0>0$. Then $\bar G$ satisfies $H_U(\eta)$ if and only if there is a probability measure $\nu$ on $[1,\infty]$ with
\[
\bar G(K_{\max}z)=\bar G_0\,z^{-\eta}\,\nu((z,\infty]),\qquad z\ge1,
\]
and then, for measurable $\phi$ with $\int_{K_{\max}}^\infty\bar G|\phi|\,du<\infty$,
\[
\int_{K_{\max}}^\infty\bar G(u)\phi(u)\,du=\bar G_0K_{\max}\int_{[1,\infty]}\int_1^{z_c}z^{-\eta}\phi(K_{\max}z)\,dz\,\nu(dz_c).
\]
The proof is that of the lemma, with two steps that are not pure substitutions. First, the boundary value is the mass at or above $K_{\max}$. Under $H_U(\eta)$, $\bar h_U(z)=z^\eta\bar G(K_{\max}z)/\bar G_0$ is nonincreasing and right-continuous on $[1,\infty)$ with $0\le\bar h_U\le\bar h_U(1)=\Q(X>K_{\max})/\bar G_0\le1$; $\nu$ is the measure with $\nu((z,\infty])=\bar h_U(z)$ for $z\ge1$, so that $\nu(\{\infty\})=\lim_{z\to\infty}\bar h_U(z)$, and with the remaining mass $1-\bar h_U(1)=\Q(X=K_{\max})/\bar G_0$ at $1$. An atom of $\nu$ at $1$ is thus an atom of $X$ at $K_{\max}$, as an atom of $\mu$ at $1$ is an atom of $X$ at $K_{\min}$. Conversely, for any such $\nu$ the displayed function is nonincreasing, right-continuous and at most $\bar G_0$, $z^\eta$ times it is nonincreasing, and it tends to $0$ as $z\to\infty$; with the atom $\bar G_0\nu(\{1\})$ placed at $K_{\max}$ it describes a law on $[K_{\max},\infty)$ of mass $\bar G_0$ that satisfies $H_U(\eta)$. Second, Tonelli's theorem is applied to $\nu((z,\infty])=\int\mathbf 1\{z<z_c\}\,\nu(dz_c)$, so the inner integral runs from $1$ to $z_c$, and the point $z_c=\infty$, which carries the pure Pareto part of the tail, contributes $\int_1^\infty z^{-\eta}\phi(K_{\max}z)\,dz$.

For $w$ integrable against $P$ on $(0,K_{\min}]$, write $W(u)=\int_u^{K_{\min}}w$, so that $\int_0^{K_{\min}}wP\,dK=\int_0^{K_{\min}}GW\,du$ by Tonelli's theorem, and define on $[0,1]$
\[
a(z_c)=\frac{1-z_c^{\gamma+1}}{\gamma+1},\qquad b_w(z_c)=\int_{z_c}^1z^\gamma W(K_{\min}z)\,dz=\frac{K_{\min}}{\gamma+1}\int_{z_c}^1w(K_{\min}z)\bigl(z^{\gamma+1}-z_c^{\gamma+1}\bigr)dz.
\]
For $w$ integrable against $C$ on $[K_{\max},\infty)$, write $W_U(u)=\int_{K_{\max}}^uw$, so that $\int_{K_{\max}}^\infty wC\,dK=\int_{K_{\max}}^\infty\bar GW_U\,du$ by Tonelli's theorem, and define on $[1,\infty]$ the upper-tail counterparts
\begin{align*}
a_U(z_c)&=\int_1^{z_c}z^{-\eta}dz=\frac{1-z_c^{1-\eta}}{\eta-1},\\
b^U_w(z_c)&=\int_1^{z_c}z^{-\eta}W_U(K_{\max}z)\,dz=\frac{K_{\max}}{\eta-1}\int_1^{z_c}w(K_{\max}z)\bigl(z^{1-\eta}-z_c^{1-\eta}\bigr)dz,
\end{align*}
with $z_c^{1-\eta}=0$ at $z_c=\infty$. By the lemma with $\phi=1$, $P_0=G_0K_{\min}\int a\,d\mu$, and with $\phi=W$, $\int_0^{K_{\min}}wP\,dK=G_0K_{\min}\int b_w\,d\mu$. By the upper-tail representation with $\phi=1$ and $\phi=W_U$, $C_0=\bar G_0K_{\max}\int a_U\,d\nu$ and $\int_{K_{\max}}^\infty wC\,dK=\bar G_0K_{\max}\int b^U_w\,d\nu$. Here $\eta>1$ is used: $a_U$ is continuous and strictly increasing on $[1,\infty]$, from $a_U(1)=0$ to its supremum $a_U(\infty)=1/(\eta-1)$, which is finite only because $\eta>1$ (for $\eta\le1$, $a_U$ is unbounded). By contrast, $a$ decreases from $a(0)=1/(\gamma+1)$ to $a(1)=0$ for every $\gamma>0$. In carrying the lower-tail arguments over, the roles of $K_{\min}$, $P_0$, $G_0$, $a$, $b_w$ and $\mu$ on $[0,1]$ are therefore played by $K_{\max}$, $C_0$, $\bar G_0$, $a_U$, $b^U_w$ and $\nu$ on $[1,\infty]$. This is a correspondence of roles, not a substitution of symbols: the normalising constant $\gamma+1$ becomes $\eta-1$, but the power $z^{\gamma+1}$ becomes $z^{1-\eta}$, the power $z^\gamma$ becomes $z^{-\eta}$, integrals over $[z_c,1]$ become integrals over $[1,z_c]$, and $W$ becomes $W_U$, so the upper-tail formulas to use are those written out above. The end points $z_c=0$ (the pure power law below $K_{\min}$) and $z_c=1$ (an atom at $K_{\min}$) correspond to $z_c=\infty$ (the pure Pareto law above $K_{\max}$) and $z_c=1$ (an atom at $K_{\max}$).

\begin{proposition}
(a) \emph{Admissible exponents.} Some law compatible with $(P_0,G_0)$ satisfies $H_L(\gamma)$ if and only if $\gamma\le\bar\gamma=K_{\min}G_0/P_0-1$, and some law compatible with $(C_0,\bar G_0)$ satisfies $H_U(\eta)$ if and only if $1<\eta\le\bar\eta=1+K_{\max}\bar G_0/C_0$. Here $\bar\gamma+1$ is the strike elasticity of the put price at $K_{\min}$ and $\bar\eta-1$ is minus that of the call price at $K_{\max}$. At $\gamma=\bar\gamma$ the lower tail is unique, $G(u)=G_0(u/K_{\min})^{\bar\gamma}$ on $(0,K_{\min}]$, so that $P(K)=P_0(K/K_{\min})^{\bar\gamma+1}$ there and every tail integral is point-identified; likewise, at $\eta=\bar\eta$, $\bar G(u)=\bar G_0(u/K_{\max})^{-\bar\eta}$ for $u\ge K_{\max}$, with no atom at $K_{\max}$, so that $C(K)=C_0(K/K_{\max})^{1-\bar\eta}$ there. The hypotheses are nested: $H_L(\gamma)$ implies $H_L(\gamma')$ for $\gamma'<\gamma$, and $H_U(\eta)$ implies $H_U(\eta')$ for $1<\eta'<\eta$.

(b) \emph{Sharp identified set.} Let $0<\gamma\le\bar\gamma$, $1<\eta\le\bar\eta$ and let $w_1,\dots,w_m$ satisfy $\int_0^{K_{\min}}|w_j|K^{\gamma+1}dK<\infty$ and $\int_{K_{\max}}^\infty|w_j|K^{1-\eta}dK<\infty$. The set of vectors $(\int_0^{K_{\min}}w_jP\,dK)_{j\le m}$ over the laws compatible with the observations and $H_L(\gamma)$ is
\[
S_L=G_0K_{\min}\Bigl\{\Bigl(\int b_{w_j}d\mu\Bigr)_{j\le m}:\ \mu\text{ a probability on }[0,1],\ \int a\,d\mu=\frac{P_0}{G_0K_{\min}}\Bigr\},
\]
and the set of vectors $(\int_{K_{\max}}^\infty w_jC\,dK)_{j\le m}$ over the laws compatible with the observations and $H_U(\eta)$ is
\[
S_U=\bar G_0K_{\max}\Bigl\{\Bigl(\int b^U_{w_j}d\nu\Bigr)_{j\le m}:\ \nu\text{ a probability on }[1,\infty],\ \int a_U\,d\nu=\frac{C_0}{\bar G_0K_{\max}}\Bigr\}.
\]
Each is a compact convex set whose extreme points are all attained by measures with at most two atoms. The two tails vary independently, so the identified set of the raw moments $(m_1,m_2,m_3)$ of R3 is the Minkowski sum $\mathrm{mid}+S_L+S_U$, with $\mathrm{mid}$ the observed middle part. For $m=1$, $S_L$ is $G_0K_{\min}$ times the interval between the lower convex and upper concave envelopes, at abscissa $P_0/(G_0K_{\min})$, of the curve $z_c\mapsto(a(z_c),b_w(z_c))$, $z_c\in[0,1]$, and $S_U$ is $\bar G_0K_{\max}$ times the corresponding interval, at abscissa $C_0/(\bar G_0K_{\max})$, for the curve $z_c\mapsto(a_U(z_c),b^U_w(z_c))$, $z_c\in[1,\infty]$.

(c) \emph{R4 revisited.} If $w\ge0$, the upper end of the identified interval of $\int_0^{K_{\min}}wP\,dK$ is R4's bound $P_0K_{\min}^{-(\gamma+1)}\int_0^{K_{\min}}wK^{\gamma+1}dK$, attained by a power law below $K_{\min}$ with an atom at $K_{\min}$, and the lower end is $G_0K_{\min}b_w(z^*)$ with $a(z^*)=P_0/(G_0K_{\min})$, attained by a law whose lower tail lies in $[K_{\min}z^*,K_{\min}]$. R4's upper bound is therefore sharp. If, in addition, $w>0$ on a set of positive measure in $(K_{\min}z^*,K_{\min})$, the lower end is positive, so R4's lower bound of zero is not sharp. In the upper tail, if $w\ge0$, the upper end of the identified interval of $\int_{K_{\max}}^\infty wC\,dK$ is R4's bound $C_0K_{\max}^{\eta-1}\int_{K_{\max}}^\infty wK^{1-\eta}dK$, attained by a Pareto law above $K_{\max}$ with an atom at $K_{\max}$, and the lower end is $\bar G_0K_{\max}b^U_w(z^*_U)$ with $a_U(z^*_U)=C_0/(\bar G_0K_{\max})$, attained by a law whose upper tail lies in $[K_{\max},K_{\max}z^*_U]$. If, in addition, $w>0$ on a set of positive measure in $(K_{\max},K_{\max}z^*_U)$, this lower end is positive and R4's lower bound of zero is not sharp in the upper tail either.

(d) \emph{Setting (i).} Suppose $G$ has a density $g$ that is continuous and positive at $K_{\min}$, and let $\gamma_i=K_{\min}g(K_{\min})/G_0$. Then $H_L(\gamma)$ implies $\gamma\le\gamma_i$, and $H_L(\gamma_i)$ holds if and only if the reversed hazard rate of $\ln X$, $ug(u)/G(u)$ at $x=\ln u$, attains its infimum over $(0,K_{\min}]$ at $K_{\min}$, which is the case if $\ln X$ has a log-concave distribution function. If $H_L(\gamma_i)$ holds, then $\gamma_i\le\bar\gamma$. With $\varepsilon(K)=KP'(K)/P(K)$ and $\gamma_i(K)$, $\bar\gamma(K)$ the same quantities at a general strike, $K\varepsilon'(K)=\varepsilon(K)\bigl(\gamma_i(K)-\bar\gamma(K)\bigr)$, so $\gamma_i\le\bar\gamma$ at $K_{\min}$ exactly when $\ln P$ is locally concave in $\ln K$ there. In the upper tail, suppose $g$ is continuous and positive at $K_{\max}$, and let $\eta_i=K_{\max}g(K_{\max})/\bar G_0$. Then $H_U(\eta)$ implies $\eta\le\eta_i$, and, for $\eta_i>1$, $H_U(\eta_i)$ holds if and only if the hazard rate of $\ln X$, $ug(u)/\bar G(u)$ at $x=\ln u$, attains its infimum over $[K_{\max},\infty)$ at $K_{\max}$, which is the case if $\ln X$ has a log-concave survival function. If $H_U(\eta_i)$ holds, then $\eta_i\le\bar\eta$. With $\varepsilon_U(K)=-KC'(K)/C(K)$ and $\eta_i(K)$, $\bar\eta(K)$ the same quantities at a general strike, $K\varepsilon_U'(K)=\varepsilon_U(K)\bigl(\bar\eta(K)-\eta_i(K)\bigr)$, so $\eta_i\le\bar\eta$ at $K_{\max}$ exactly when $\ln C$ is locally concave in $\ln K$ there, in the same sense.

(e) \emph{Breakdown.} For $\theta\in(0,1]$ let $I(\theta)$ be the identified set of the variance and third central moment under $H_L(\theta\bar\gamma)$ and $H_U(1+\theta(\bar\eta-1))$. Then $I(\theta')\subseteq I(\theta)$ for $\theta<\theta'$, and $I(1)$ is a single point. Say that the sign of the third central moment, and hence of the skewness, is identified at $\theta$ if the third central moment is positive at every point of $I(\theta)$ or negative at every point of $I(\theta)$. Because the sets are nested, the set of $\theta$ at which the sign is identified is empty or an interval with right end $1$. It is empty exactly when the third central moment at the single point of $I(1)$ is zero; there is then no breakdown point, and the sign is identified at no fraction. Otherwise its infimum $\theta^*\in[0,1]$ is the breakdown point of the sign; that $\theta^*>0$ is not shown here.
\end{proposition}

\begin{proof}
(a) By the lemma, $P_0/(G_0K_{\min})=\int a\,d\mu$, and $a$ is continuous and strictly decreasing from $a(0)=1/(\gamma+1)$ to $a(1)=0$. A feasible $\mu$ therefore exists if and only if $P_0/(G_0K_{\min})\le1/(\gamma+1)$, that is $\gamma\le\bar\gamma$; for such $\gamma$ the mixture of $\delta_0$ and $\delta_1$ with weight $(\gamma+1)P_0/(G_0K_{\min})$ on $\delta_0$ is feasible. At $\gamma=\bar\gamma$ the constraint reads $\int a\,d\mu=a(0)$, which forces $\mu=\delta_0$ because $a<a(0)$ on $(0,1]$. In the upper tail, $C_0/(\bar G_0K_{\max})=\int a_U\,d\nu$ by the upper-tail representation, and $a_U$ is continuous and strictly increasing from $a_U(1)=0$ to its supremum $a_U(\infty)=1/(\eta-1)$, finite because $\eta>1$. A feasible $\nu$ therefore exists if and only if $C_0/(\bar G_0K_{\max})\le1/(\eta-1)$, that is $\eta\le\bar\eta$; for such $\eta>1$ the mixture of $\delta_1$ and $\delta_\infty$ with weight $(\eta-1)C_0/(\bar G_0K_{\max})$ on $\delta_\infty$ is feasible. At $\eta=\bar\eta$ the constraint reads $\int a_U\,d\nu=a_U(\infty)$, which forces $\nu=\delta_\infty$ because $a_U<a_U(\infty)$ on $[1,\infty)$; the representation then gives $\bar G(u)=\bar G_0(u/K_{\max})^{-\bar\eta}$ for $u\ge K_{\max}$ and, integrating, $C(K)=\bar G_0K_{\max}(K/K_{\max})^{1-\bar\eta}/(\bar\eta-1)=C_0(K/K_{\max})^{1-\bar\eta}$. The lower limit $\eta>1$ is part of $H_U(\eta)$, and $\bar\eta>1$ always, since $C_0,\bar G_0>0$. The elasticity statements are the definitions, since $P'(K_{\min}+)=G_0$ and $C'(K_{\max}-)=-\bar G_0$. Nesting holds because $G/u^{\gamma'}=(G/u^\gamma)\,u^{\gamma-\gamma'}$ is a product of nonnegative nondecreasing functions, and, for $1<\eta'<\eta$, $u^{\eta'}\bar G=(u^\eta\bar G)\,u^{\eta'-\eta}$ is a product of nonnegative nonincreasing functions.

(b) By the lemma the map $\mu\mapsto$ law is a bijection between the probability measures on $[0,1]$ satisfying the constraint and the laws below $K_{\min}$ compatible with $(P_0,G_0)$ and $H_L(\gamma)$; each such law extends to a law on $(0,\infty)$ compatible with the observed prices, because the law on $(K_{\min},K_{\max})$ is fixed by the smile and the mass and the mean $F$ are fixed by $P_0$, $G_0$, $C_0$ and $\bar G_0$ (R2). The integrability conditions make $a$ and $b_{w_j}$ continuous on $[0,1]$, by dominated convergence with the bound $|w_j(K_{\min}z)|z^{\gamma+1}$, so the feasible set of $\mu$ is weak$^*$ compact and convex and its image under the continuous linear map $\mu\mapsto(\int b_{w_j}d\mu)_j$ is compact and convex. Extreme points of the image are images of extreme points of the feasible set, which are mixtures of at most two Dirac measures (Winkler, 1988, Theorem 2.1, as quoted by Pinelis, 2012). For $m=1$ the set of pairs $(\int a\,d\mu,\int b_w\,d\mu)$ over all probability measures is the convex hull of the compact curve $\{(a(z_c),b_w(z_c))\}$ (Kemperman, 1968), and its section at abscissa $P_0/(G_0K_{\min})$ is the stated interval. The upper tail is the same argument on $[1,\infty]$, which is compact (the map $z\mapsto1/z$ takes it onto $[0,1]$), with the upper-tail representation in place of the lemma. The only step that is not a pure substitution is continuity at $z_c=\infty$: $a_U$ is continuous there because $\eta>1$, and $b^U_{w_j}$ is continuous on $[1,\infty]$ by dominated convergence with the bound $|w_j(K_{\max}z)|z^{1-\eta}$, which is integrable on $[1,\infty)$ by the integrability condition. The lower-tail law and the upper-tail law can be chosen separately, which gives the Minkowski sum.

(c) For $w\ge0$, $W$ is nonnegative and nonincreasing. Along the curve, $db_w/dz_c=-z_c^\gamma W(K_{\min}z_c)$ and $da/dz_c=-z_c^\gamma$, so $b_w$ is a function of $a$ with slope $W(K_{\min}z_c)$, which increases as $a$ increases. The curve is therefore the graph of a convex function of $a$: its lower envelope is the curve itself, attained by the single atom $z^*$, and its upper envelope is the chord from $(0,0)$ at $z_c=1$ to $(a(0),b_w(0))$ at $z_c=0$. At abscissa $P_0/(G_0K_{\min})$ the chord gives $G_0K_{\min}\cdot(\gamma+1)\frac{P_0}{G_0K_{\min}}\,b_w(0)=(\gamma+1)P_0\int_0^{K_{\min}}w(K)\frac{(K/K_{\min})^{\gamma+1}}{\gamma+1}dK$, which is R4's bound. The law is the mixture of $\delta_0$ and $\delta_1$ of the proof of (a). If $w>0$ on a set of positive measure in $(K_{\min}z^*,K_{\min})$, then $b_w(z^*)>0$, because $z^{\gamma+1}-z^{*\gamma+1}>0$ for $z>z^*$; here $z^*<1$ because $P_0>0$.

In the upper tail, for $w\ge0$, $W_U$ is nonnegative and nondecreasing. Along the curve, $db^U_w/dz_c=z_c^{-\eta}W_U(K_{\max}z_c)$ and $da_U/dz_c=z_c^{-\eta}$, so $b^U_w$ is a function of $a_U$ with slope $W_U(K_{\max}z_c)$, which increases as $a_U$ increases, since both increase with $z_c$. The curve is therefore again the graph of a convex function, now on $[0,1/(\eta-1)]$: its lower envelope is the curve itself, attained by the single atom $z^*_U$, which exists and is unique because $a_U$ is a continuous strictly increasing map of $[1,\infty]$ onto $[0,1/(\eta-1)]$ and $C_0/(\bar G_0K_{\max})$ lies in that interval by (a); its upper envelope is the chord from $(0,0)$ at $z_c=1$ to $(1/(\eta-1),b^U_w(\infty))$ at $z_c=\infty$. At abscissa $C_0/(\bar G_0K_{\max})$ the chord gives
\begin{align*}
\bar G_0K_{\max}\cdot(\eta-1)\frac{C_0}{\bar G_0K_{\max}}\,b^U_w(\infty)&=C_0K_{\max}\int_1^\infty w(K_{\max}z)\,z^{1-\eta}dz\\
&=C_0K_{\max}^{\eta-1}\int_{K_{\max}}^\infty w(K)K^{1-\eta}dK,
\end{align*}
which is R4's bound. The law is the mixture of $\delta_1$ and $\delta_\infty$ of the proof of (a): an atom of mass $\bar G_0-(\eta-1)C_0/K_{\max}$ at $K_{\max}$ and $\bar G(u)=(\eta-1)(C_0/K_{\max})(u/K_{\max})^{-\eta}$ for $u\ge K_{\max}$. The single atom $\nu=\delta_{z^*_U}$ gives $\bar G(K_{\max}z)=\bar G_0z^{-\eta}$ for $1\le z<z^*_U$ and $0$ beyond, so the upper tail lies in $[K_{\max},K_{\max}z^*_U]$. If $w>0$ on a set of positive measure in $(K_{\max},K_{\max}z^*_U)$, then $b^U_w(z^*_U)>0$, because $z^{1-\eta}-z_U^{*\,1-\eta}>0$ for $z<z^*_U$; here $z^*_U>1$ because $C_0>0$.

(d) $\ln(G/u^\gamma)$ is nondecreasing on $(0,K_{\min}]$ and differentiable from the left at $K_{\min}$ with derivative $g/G-\gamma/u$, which is therefore nonnegative: $\gamma\le\gamma_i$. With a density, $\ln G$ is locally absolutely continuous where $G>0$, so $H_L(\gamma)$ holds if and only if $ug(u)/G(u)\ge\gamma$ almost everywhere on $(0,K_{\min}]$; with $\gamma=\gamma_i$ this says that the infimum is attained at $K_{\min}$. If $\ln X$ has a log-concave distribution function, its reversed hazard rate is nonincreasing and the infimum over $x\le\ln K_{\min}$ is attained at the end point; a log-concave density suffices by Bagnoli and Bergstrom (2005, Theorem 1). The inequality $\gamma_i\le\bar\gamma$ is (a) applied to $\gamma_i$. Finally, $\varepsilon=KG/P$ gives $K\varepsilon'=\varepsilon+\varepsilon Kg/G-\varepsilon^2=\varepsilon(\gamma_i(K)+1-\varepsilon)=\varepsilon(\gamma_i(K)-\bar\gamma(K))$, and $K\varepsilon'(K)$ is the derivative of $\varepsilon$ with respect to $\ln K$.

In the upper tail, a density continuous at $K_{\max}$ excludes an atom there, so $\bar G(K_{\max})=\bar G_0$. $\ln(u^\eta\bar G)$ is nonincreasing on $[K_{\max},\infty)$ and differentiable from the right at $K_{\max}$ with derivative $\eta/u-g/\bar G$, which is therefore nonpositive: $\eta\le\eta_i$. With a density, $\ln\bar G$ is locally absolutely continuous where $\bar G>0$, so $H_U(\eta)$ holds if and only if $ug(u)/\bar G(u)\ge\eta$ almost everywhere on $[K_{\max},\infty)$ where $\bar G>0$; at $x=\ln u$ this ratio is the hazard rate of $\ln X$, and with $\eta=\eta_i$ the condition says that its infimum is attained at $K_{\max}$. If $\ln X$ has a log-concave survival function, its hazard rate is nondecreasing and the infimum over $x\ge\ln K_{\max}$ is attained at the end point; a log-concave density of $\ln X$ suffices, by Bagnoli and Bergstrom (2005, Theorem 1) applied to $-\ln X$, whose distribution function is $x\mapsto\Q(\ln X\ge-x)$. The inequality $\eta_i\le\bar\eta$ is (a) applied to $\eta_i$. Finally, $\varepsilon_U=K\bar G/C$ and $C'=-\bar G$ give $K\varepsilon_U'=\varepsilon_U-\varepsilon_UKg/\bar G+\varepsilon_U^2=\varepsilon_U(1+\varepsilon_U-\eta_i(K))=\varepsilon_U(\bar\eta(K)-\eta_i(K))$, and $-\varepsilon_U$ is the derivative of $\ln C$ with respect to $\ln K$, so $\eta_i\le\bar\eta$ exactly when $d^2\ln C/d(\ln K)^2\le0$ at $K_{\max}$.

(e) By the nesting in (a), the laws compatible with the hypotheses at $\theta'$ are compatible with them at $\theta$, so the identified sets are nested; at $\theta=1$ both tails are unique by (a), so $I(1)$ is a single point, and every $I(\theta)$ contains it. If the sign is identified at $\theta$, it is identified at every $\theta'>\theta$, because $I(\theta')\subseteq I(\theta)$; the set is therefore empty or an interval with right end $1$. If the third central moment at the point of $I(1)$ is nonzero, the sign is identified at $\theta=1$ and the set is not empty. If it is zero, every $I(\theta)$ contains a point at which the third central moment is zero, so the sign is identified at no $\theta$.
\end{proof}

\begin{remark}
The implementation (\texttt{src/qef/fx/identification.py}) restricts $\mu$ and $\nu$ to a grid of step points graded towards the boundary, which gives an inner approximation that converges as the grid is refined. For a fixed mean, the ranges of the variance and of the third central moment are computed exactly for that grid by linear programmes in the measures. Their extremes over the mean are not computed exactly: \texttt{SharpMomentSet.\_extreme} evaluates them at 25 means spread over the range of the mean (15 in the bisection for the breakdown point of (e)) and refines around the two best by bounded scalar minimisation. Every mean it evaluates is the mean of a feasible law, and every value it returns is attained by such a law, so the computed ranges are inner approximations of the sharp set for the step grid; they can only err towards declaring the sign of the skewness identified. A check with eight times as many means on 40 currency-months is reported in \texttt{reports/stage4.md}. The weights of the moment contracts of R3 satisfy the integrability conditions of (b) for every $\gamma>0$ and $\eta>1$. Because $I(1)$ is a point, the Pareto tails of (a) give the unique moments at the strongest admissible hypothesis; the breakdown point, when it exists, measures how far below that hypothesis the sign of the skewness survives. In the model validation of \texttt{reports/model\_validation.md}, the model's own tail exponents lie between $0.58$ and $0.97$ of the maximal ones, setting (i) fails in seven of eight cases, and the identified set at the model's own exponents contains the true moments in every case, as (b) requires.
\end{remark}

An observed price beyond a boundary strike tests the hypotheses and, if consistent, adds a linear constraint on the mixing measure. The following consequence of the lemma is used for the out-of-sample test with $5\Delta$ quotes.

\begin{proposition}
Let $0<K<K_{\min}$, $k=K/K_{\min}$ and $0\le\gamma\le\bar\gamma$, where $\gamma=0$ stands for the limit of the weakest hypothesis. Under $H_L(\gamma)$ the identified interval of $P(K)$ is
\[
\Bigl[\max\Bigl(0,\;P_0-G_0K_{\min}\,\frac{1-k^{\gamma+1}}{\gamma+1}\Bigr),\;P_0k^{\gamma+1}\Bigr].
\]
For $K>K_{\max}$, $k=K/K_{\max}$ and $1\le\eta\le\bar\eta$, the identified interval of $C(K)$ under $H_U(\eta)$ is $[\max(0,C_0-\bar G_0K_{\max}a^U_k),\,C_0k^{1-\eta}]$ with $a^U_k=(1-k^{1-\eta})/(\eta-1)$, read as $\ln k$ at $\eta=1$, where $\eta=1$ again stands for the limit of the weakest hypothesis. At the limits $\gamma=0$ and $\eta=1$ the upper ends are approached but not attained. The upper ends decrease and the lower ends increase with the exponent, so an observed price at $K$ is consistent with the hypothesis exactly for exponents in an interval that starts at the weakest hypothesis and may be empty.
\end{proposition}

\begin{proof}
By the lemma with $\phi=\mathbf 1\{u\le K\}$, $P(K)=\int_0^{K_{\min}}G\mathbf 1\{u\le K\}du=G_0K_{\min}\int b\,d\mu$ with $b(z_c)=\int_{z_c}^1z^\gamma\mathbf 1\{z\le k\}dz=\max(0,a(z_c)-a_k)$, where $a_k=(1-k^{\gamma+1})/(\gamma+1)$. As a function of $a$ this is convex and piecewise linear, so the lower envelope at $P_0/(G_0K_{\min})$ is the function itself, which gives the lower end, and the upper envelope is the chord from $(0,0)$ to $(a(0),b(0))=(1/(\gamma+1),k^{\gamma+1}/(\gamma+1))$, which gives $P_0k^{\gamma+1}$, R4's bound. Since $a_k=\int_k^1z^\gamma dz$ decreases in $\gamma$ for $k<1$, the lower end increases in $\gamma$, and $P_0k^{\gamma+1}$ decreases. In the upper tail, the upper-tail representation with $\phi=\mathbf 1\{u\ge K\}$ gives $C(K)=\int_{K_{\max}}^\infty\bar G\,\mathbf 1\{u\ge K\}du=\bar G_0K_{\max}\int b^U\,d\nu$ with $b^U(z_c)=\int_1^{z_c}z^{-\eta}\mathbf 1\{z\ge k\}dz=\max(0,a_U(z_c)-a^U_k)$, where $a^U_k=a_U(k)=\int_1^kz^{-\eta}dz$. For $\eta>1$ this is again convex and piecewise linear in $a_U$, so the lower envelope at $C_0/(\bar G_0K_{\max})$ gives the lower end, and the upper envelope is the chord from $(0,0)$ at $z_c=1$ to $(a_U(\infty),b^U(\infty))=(1/(\eta-1),k^{1-\eta}/(\eta-1))$ at $z_c=\infty$, which gives $\bar G_0K_{\max}\cdot(\eta-1)\frac{C_0}{\bar G_0K_{\max}}\cdot\frac{k^{1-\eta}}{\eta-1}=C_0k^{1-\eta}$, R4's bound. Since $a^U_k$ decreases in $\eta$ for $k>1$, the lower end increases in $\eta$, and $C_0k^{1-\eta}$ decreases. At $\eta=1$, $a_U(z_c)=\ln z_c$ is unbounded and the chord has no right end point. The lower end is unchanged, with $a^U_k=\ln k$, while the upper end $C_0=\lim_{\eta\downarrow1}C_0k^{1-\eta}$ is approached by mixtures of $\delta_1$ and $\delta_z$ as $z\to\infty$ but not attained: $C(K)=C_0$ would require $\bar G=0$ on $(K_{\max},K)$, hence $C_0=0$.

The lower tail at $\gamma=0$ is similar. This case lies outside the lemma's hypothesis $\gamma>0$: $H_L(0)$ says only that $G$ is nondecreasing, and the representation of the lemma still holds with $z^0=1$, except that the displayed function now tends to $G_0\mu(\{0\})$ as $z\downarrow0$, so an atom of $\mu$ at $0$ would be an atom of $X$ at $0$, which R3 excludes ($X>0$). The feasible $\mu$ therefore have no atom at $0$, while the chord to $z_c=0$ that gives the upper end needs one. The lower end is unchanged, and it is attained when $\bar\gamma>0$, by the single atom at $z^*=\bar\gamma/(\bar\gamma+1)>0$. The upper end $P_0k=\lim_{\gamma\downarrow0}P_0k^{\gamma+1}$ is approached by mixtures of $\delta_1$ and $\delta_z$ as $z\downarrow0$ but not attained, as $C_0$ is not attained at $\eta=1$: since $G$ is nondecreasing, $P(K)=P_0k$, that is, equal averages of $G$ over $(0,K)$ and $(0,K_{\min})$, would require $G$ to be constant almost everywhere on $(0,K_{\min})$, hence equal to $G(0+)=\Q(X\le0)=0$, and then $P_0=0$.
\end{proof}

\begin{remark}
In the out-of-sample test (\texttt{scripts/estimate\_wing\_test.py}) the $5\Delta$ put and call prices also enter the identified set as equality constraints, one on each mixing measure, with the row $\max(0,a(z_c)-a_k)$ in the lower tail and its counterpart in the upper tail, and the breakdown point is searched along the common path $\gamma=\theta\bar\gamma$, $\eta=1+\theta(\bar\eta-1)$. Write $I_5(\theta)$ for the constrained identified set, and $\theta_5=\min(\theta_5^L,\theta_5^U)$, where $\theta_5^L$ and $\theta_5^U$ are the largest fractions of the maximal exponents at which the $5\Delta$ put and call prices lie in their intervals above. Two parts of R10 do not carry over. First, (e) is proved without extra constraints, where $I(1)$ is a single point. The constraints do not depend on $\theta$, so the sets $I_5(\theta)$ are still nested, and they are not empty exactly for $\theta\le\theta_5$; the fractions in $(0,\theta_5]$ at which the sign is identified are therefore empty or an interval with right end $\theta_5$. Above $\theta_5$ there is no compatible law, and $I_5(\theta_5)$ is in general not a single point, so the interval structure of (e), with right end $1$ and its criterion for when the set is empty, is not proved in this case. Second, each mixing measure is now subject to two linear constraints besides its total mass, so by Winkler's theorem the extreme points of its feasible set are mixtures of at most three Dirac measures, and three atoms occur: if $\gamma<\bar\gamma$ and the $5\Delta$ put price lies strictly inside its interval, the point $\bigl(P_0,P(K)\bigr)/(G_0K_{\min})$ lies in the interior of the triangle with vertices $(a(z_c),\max(0,a(z_c)-a_k))$ at $z_c=0,k,1$, and the mixture of $\delta_0$, $\delta_k$ and $\delta_1$ that matches both prices is an extreme point, because the vectors $(1,a(z_c),\max(0,a(z_c)-a_k))$ at these three points are linearly independent. The tail sets remain compact and convex, but the two-atom description of (b) does not hold. With the constraints, the breakdown point $\theta^*_5$ is therefore defined computationally, as the script computes it. Let $\theta_{\max}=0.999\,\theta_5$. If $\theta_{\max}\le0.01$ or the sign is not identified at $\theta_{\max}$, $\theta^*_5$ is undefined; if the sign is identified at $\theta=0.01$, $\theta^*_5=0.01$; otherwise $\theta^*_5$ is the smallest fraction found by bisection on $[0.01,\theta_{\max}]$, with tolerance $0.01$, at which the sign is identified, each evaluation being computed as in the remark after the proof of R10, with 25 means and the two extra constraints. The bisection relies on the nesting, which holds for the sharp sets on $(0,\theta_5]$.
\end{remark}

\section*{The Garman--Kohlhagen PDE}

The equation, and the representation of its solution as an expectation, are standard (Garman and Kohlhagen, 1983; for the Feynman--Kac argument, Shreve, 2004). The verification and the uniqueness argument below are written out for completeness.

Let $dS_t=(r_d-r_f)S_t\,dt+\sigma S_t\,dW_t$ under $\Q^d$ with constant coefficients, and let $c(t,S)$ be the Garman and Kohlhagen (1983) call price with strike $K$ and expiry $T$,
\[
c(t,S)=Se^{-r_f(T-t)}\Phi(d_+)-Ke^{-r_d(T-t)}\Phi(d_-),\qquad d_\pm=\frac{\ln(S/K)+(r_d-r_f\pm\sigma^2/2)(T-t)}{\sigma\sqrt{T-t}}.
\]

\begin{proposition}
(a) $c\in C^{1,2}([0,T)\times(0,\infty))$ satisfies
\[
\partial_tc+(r_d-r_f)S\,\partial_Sc+\tfrac12\sigma^2S^2\,\partial_{SS}c-r_dc=0,
\]
$0\le c\le Se^{-r_f(T-t)}$, and $c(t,S)\to(S-K)^+$ as $t\uparrow T$, locally uniformly in $S$. (b) If $u\in C^{1,2}([0,T)\times(0,\infty))\cap C([0,T]\times(0,\infty))$ solves the same equation with $u(T,S)=(S-K)^+$ and $|u(t,S)|\le a(1+S^m)$ for some $a,m$, then $u=c$. No boundary condition at $S=0$ is required.
\end{proposition}

\begin{proof}
(a) Write $\tau=T-t$. The identity $Se^{-r_f\tau}\varphi(d_+)=Ke^{-r_d\tau}\varphi(d_-)$ gives $\partial_Sc=e^{-r_f\tau}\Phi(d_+)$ and $\partial_{SS}c=e^{-r_f\tau}\varphi(d_+)/(S\sigma\sqrt\tau)$. The same identity gives $\partial_tc=r_fSe^{-r_f\tau}\Phi(d_+)-r_dKe^{-r_d\tau}\Phi(d_-)-Se^{-r_f\tau}\varphi(d_+)\sigma/(2\sqrt\tau)$. Substitution verifies the equation. The bounds follow from $c=e^{-r_d\tau}\E[(S_T-K)^+\mid S_t=S]$ and $\E[S_T\mid S_t=S]=Se^{(r_d-r_f)\tau}$. As $\tau\downarrow0$, $\Phi(d_\pm)\to\mathbf 1\{S>K\}$ for $S\ne K$ and $c(t,K)\to0$, and $c$ is Lipschitz in $S$ with constant $e^{|r_f|T}$ uniformly in $t$, which gives locally uniform convergence.

(b) Fix $(t,S)$ and let $\tau_n=\inf\{s\ge t:S_s\notin(1/n,n)\}\wedge(T-1/n)$. It\^o's formula for $e^{-r_d(s-t)}u(s,S_s)$ on $[t,\tau_n]$, the equation, and the boundedness of $\partial_Su$ on compact sets give $u(t,S)=\E\bigl[e^{-r_d(\tau_n-t)}u(\tau_n,S_{\tau_n})\bigr]$. Because $S$ is a geometric Brownian motion started at $S>0$, it reaches neither $0$ nor $\infty$ before $T$, so $\tau_n\to T$ almost surely. Continuity of $u$ on $[0,T]\times(0,\infty)$ then gives $u(\tau_n,S_{\tau_n})\to(S_T-K)^+$. The growth bound and $\E\sup_{s\le T}S_s^m<\infty$ allow dominated convergence, so $u(t,S)=e^{-r_d(T-t)}\E(S_T-K)^+=c(t,S)$. The representation uses only values in the interior $(0,\infty)$, which the process never leaves, so no condition at $S=0$ enters. It\^o's formula and the Feynman--Kac representation are standard results, stated for example in Shreve (2004).
\end{proof}

\section*{References}
\begin{list}{}{\leftmargin=1.5em \itemindent=-1.5em \itemsep=2pt \parsep=0pt}\small
\item Alvarez, F. and Jermann, U. J. (2005). Using asset prices to measure the persistence of the marginal utility of wealth. \emph{Econometrica} 73(6), 1977--2016. Not consulted directly; cited as stated in Backus, Chernov and Zin (2014).
\item Backus, D. K., Foresi, S. and Telmer, C. I. (2001). Affine term structure models and the forward premium anomaly. \emph{Journal of Finance} 56(1), 279--304. Versions consulted: published article (JSTOR copy), Sections I.A--I.B, and NBER Working Paper 5623 (1996).
\item Backus, D., Chernov, M. and Zin, S. (2014). Sources of entropy in representative agent models. \emph{Journal of Finance} 69(1), 51--99.
\item Bagnoli, M. and Bergstrom, T. (2005). Log-concave probability and its applications. \emph{Economic Theory} 26(2), 445--469.
\item Bakshi, G., Kapadia, N. and Madan, D. (2003). Stock return characteristics, skew laws, and the differential pricing of individual equity options. \emph{Review of Financial Studies} 16(1), 101--143.
\item Bansal, R. and Lehmann, B. N. (1997). Growth-optimal portfolio restrictions on asset pricing models. \emph{Macroeconomic Dynamics} 1(2), 333--354. Not consulted directly; cited as stated in Backus, Chernov and Zin (2014).
\item Benaim, S., Dodgson, M. and Kainth, D. (2008). An arbitrage-free method for smile extrapolation. Technical report, Royal Bank of Scotland. Version consulted: RiskTailsPaper\_v5 of January 2009, from an Internet Archive snapshot.
\item Bollerslev, T., Todorov, V. and Xu, L. (2015). Tail risk premia and return predictability. \emph{Journal of Financial Economics} 118(1), 113--134.
\item Brandt, M. W., Cochrane, J. H. and Santa-Clara, P. (2006). International risk sharing is better than you think, or exchange rates are too smooth. \emph{Journal of Monetary Economics} 53(4), 671--698.
\item Breeden, D. T. and Litzenberger, R. H. (1978). Prices of state-contingent claims implicit in option prices. \emph{Journal of Business} 51(4), 621--651.
\item Carr, P. and Madan, D. (2001). Towards a theory of volatility trading. In E. Jouini, J. Cvitani\'c and M. Musiela (eds), \emph{Option Pricing, Interest Rates and Risk Management}, Cambridge University Press, 458--476.
\item Carr, P., Fisher, T. and Ruf, J. (2014). On the hedging of options on exploding exchange rates. \emph{Finance and Stochastics} 18(1), 115--144.
\item Choi, J. and Wu, L. (2021). The equivalent constant-elasticity-of-variance (CEV) volatility of the stochastic-alpha-beta-rho (SABR) model. \emph{Journal of Economic Dynamics and Control} 128, 104143. Version consulted: arXiv:1911.13123v5.
\item Du, W., Tepper, A. and Verdelhan, A. (2018). Deviations from covered interest rate parity. \emph{Journal of Finance} 73(3), 915--957. Version consulted: NBER Working Paper 23170, February 2017.
\item Farhi, E., Fraiberger, S. P., Gabaix, X., Ranci\`ere, R. and Verdelhan, A. (2015). Crash risk in currency markets. Working paper, version of 12 March 2015, SSRN 1397668; first issued as NBER Working Paper 15062 (2009).
\item F\"ollmer, H. and Schied, A. (2004). \emph{Stochastic Finance: An Introduction in Discrete Time}, 2nd edition. De Gruyter. Lemma 7.23 consulted in a page preview.
\item Fukasawa, M. (2012). The normalizing transformation of the implied volatility smile. \emph{Mathematical Finance} 22(4), 753--762. Version consulted: arXiv:1008.5055v2.
\item Garman, M. B. and Kohlhagen, S. W. (1983). Foreign currency option values. \emph{Journal of International Money and Finance} 2(3), 231--237.
\item Gatheral, J. and Jacquier, A. (2014). Arbitrage-free SVI volatility surfaces. \emph{Quantitative Finance} 14(1), 59--71.
\item Geman, H., El Karoui, N. and Rochet, J.-C. (1995). Changes of num\'eraire, changes of probability measure and option pricing. \emph{Journal of Applied Probability} 32(2), 443--458.
\item Hagan, P. S., Kumar, D., Lesniewski, A. S. and Woodward, D. E. (2002). Managing smile risk. \emph{Wilmott Magazine}, 84--108.
\item Henrard, M. (2011). Adjoint algorithmic differentiation: calibration and implicit function theorem. OpenGamma Quantitative Research; SSRN 1896329 (published in the \emph{Journal of Computational Finance}, 2014). Version consulted: SSRN version of 1 November 2011.
\item Horowitz, J. L. and Manski, C. F. (1995). Identification and robustness with contaminated and corrupted data. \emph{Econometrica} 63(2), 281--302. Not consulted; cited as credited by Masten and Poirier (2020).
\item J\"ackel, P. (2020). Strike from volatility and delta-with-premium. \emph{Quantitative Finance} 20(8), 1227--1235.
\item Jiang, G. J. and Tian, Y. S. (2005). The model-free implied volatility and its information content. \emph{Review of Financial Studies} 18(4), 1305--1342. Consulted: published article, Section 1.2.1 (Proposition 2).
\item Jiang, Z., Krishnamurthy, A. and Lustig, H. (2021). Foreign safe asset demand and the dollar exchange rate. \emph{Journal of Finance} 76(3), 1049--1089. Version consulted: NBER Working Paper 24439.
\item Jurek, J. W. (2014). Crash-neutral currency carry trades. \emph{Journal of Financial Economics} 113(3), 325--347. Version consulted for Section 2.2: SSRN 1262934, August 2013.
\item Kemperman, J. H. B. (1968). The general moment problem, a geometric approach. \emph{Annals of Mathematical Statistics} 39(1), 93--122.
\item Kline, P. and Santos, A. (2013). Sensitivity to missing data assumptions: theory and an evaluation of the U.S. wage structure. \emph{Quantitative Economics} 4(2), 231--267.
\item Lando, T., Arab, I. and Oliveira, P. E. (2022). Transform orders and stochastic monotonicity of statistical functionals. Working paper, arXiv:2112.02383v2, 24 October 2022.
\item Le Floc'h, F. and Kennedy, G. (2014). Explicit SABR calibration through simple expansions. Working paper, SSRN 2467231, version 1.2 of July 2014.
\item Lee, R. W. (2004). The moment formula for implied volatility at extreme strikes. \emph{Mathematical Finance} 14(3), 469--480. Version consulted: preprint of 11 August 2003.
\item Lustig, H. and Verdelhan, A. (2019). Does incomplete spanning in international financial markets help to explain exchange rates? \emph{American Economic Review} 109(6), 2208--2244. Version consulted: manuscript of January 2018.
\item Masten, M. A. and Poirier, A. (2020). Inference on breakdown frontiers. \emph{Quantitative Economics} 11(1), 41--111.
\item Milnor, J. (1963). \emph{Morse Theory}. Annals of Mathematics Studies 51. Princeton University Press.
\item Mingone, A. (2023). Smiles in delta. \emph{Quantitative Finance} 23(12), 1713--1728. Version consulted: arXiv:2209.00406v1.
\item Pinelis, I. (2012). On the extreme points of moments sets. Working paper, arXiv:1204.0249v1, 1 April 2012. Consulted for its statement of Winkler (1988, Theorem 2.1).
\item Popescu, I. (2005). A semidefinite programming approach to optimal-moment bounds for convex classes of distributions. \emph{Mathematics of Operations Research} 30(3), 632--657. Version consulted: INSEAD Working Paper 2002/15/TM.
\item Qin, H., Yang, R., Che, C. and Feng, L. (2026). Arbitrage-free multi-maturity risk-neutral marginals. Working paper, arXiv:2607.06204v1, 7 July 2026.
\item Reiswich, D. (2010). \emph{The Foreign Exchange Volatility Surface}. Doctoral dissertation, Frankfurt School of Finance \& Management.
\item Reiswich, D. and Wystup, U. (2012). FX volatility smile construction. \emph{Wilmott} 2012(60), 58--69. Version consulted: CPQF Working Paper No.~20, version 2, Frankfurt School of Finance \& Management, 20 March 2010.
\item Sampford, M. R. (1953). Some inequalities on Mill's ratio and related functions. \emph{Annals of Mathematical Statistics} 24(1), 130--132.
\item Shreve, S. E. (2004). \emph{Stochastic Calculus for Finance II: Continuous-Time Models}. Springer. Consulted: bibliographic record; cited for standard textbook statements.
\item Stoye, J. (2005, 2010). Cited as credited by Masten and Poirier (2020); not consulted.
\item S\"uli, E. and Mayers, D. F. (2003). \emph{An Introduction to Numerical Analysis}. Cambridge University Press.
\item Taleb, N. N., Yarckin, B., Mann, C., Delic, D. and Spitznagel, M. (2023). Tail option pricing under power laws. Working paper, arXiv:1908.02347v3, March 2023.
\item Veres-Ferrer, E.-J. and Pav\'ia, J. M. (2022). The elasticity of a random variable as a tool for measuring and assessing risks. \emph{Risks} 10(3), 68.
\item Winkler, G. (1988). Extreme points of moment sets. \emph{Mathematics of Operations Research} 13(4), 581--587. Not consulted directly; Theorem 2.1 as quoted by Pinelis (2012, Theorem 12).
\item Zhuang, J. (2026). The risk-neutral crash frontier: sharp joint bounds on crash probability and conditional depth from option bid-ask quotes. Working paper, arXiv:2607.25353v4, 18 September 2026.
\end{list}

\end{document}
